% 93-20(93쪽) — 두 곡선 $y=x^2-15$, $y=-x^2+15$ 로 둘러싸인 부분에 내접하고 한 쌍의 대변이 $x$축에 평행한 직사각형. 스캔 화질이 낮아 TikZ 로 다시 그림.
% 원문에 있는 것만 옮긴다: 두 포물선 · 색칠한 직사각형 · 라벨 15, $-15$, O, x, y, $y=x^2-15$, $y=-x^2+15$. 직사각형의 가로·세로 길이와 넓이는 원문에 없고 답이 읽히면 안 되므로 적지 않았고, 최댓값이 되는 자리가 아닌 일반적인 위치로 그렸다.
\begin{tikzpicture}[x=0.29cm, y=0.149cm, line width=0.6pt]
  \fill[black!12] (-2.5,-8.75) rectangle (2.5,8.75);
  \draw[->, >=stealth, line width=0.7pt] (-7,0) -- (9.1,0) node[below=2pt, font=\small] {$x$};
  \draw[->, >=stealth, line width=0.7pt] (0,-21) -- (0,23.5) node[left=1pt, font=\small] {$y$};
  \draw[line width=0.45pt] (-2.5,-8.75) rectangle (2.5,8.75);
  \draw[domain=-5.8:5.8, samples=120, smooth] plot (\x, {(\x)*(\x)-15});
  \draw[domain=-5.8:5.8, samples=120, smooth] plot (\x, {15-(\x)*(\x)});
  \node[anchor=east, font=\small] at (-0.3,15.8) {$15$};
  \node[anchor=east, font=\small] at (-0.3,-15.8) {$-15$};
  \node[below left=-1pt and -3pt, font=\small] at (0,0) {$\pt{O}$};
  \node[anchor=west, font=\small] at (1.3,19.8) {$y=x^2-15$};
  \node[anchor=west, font=\small] at (0.9,-19.8) {$y=-x^2+15$};
\end{tikzpicture}
